%-------------------------
% Resume in Latex
% Based on the Jake Gutierrez template: https://github.com/jakegut/resume
% Filled in with Amangeldy Shyngyskhan's data — RU version
%------------------------

\documentclass[letterpaper,11pt]{article}

\usepackage{latexsym}
\usepackage[empty]{fullpage}
\usepackage{titlesec}
\usepackage{marvosym}
\usepackage[usenames,dvipsnames]{color}
\usepackage{verbatim}
\usepackage{enumitem}
\usepackage[hidelinks]{hyperref}
\usepackage{fancyhdr}
\usepackage[T2A]{fontenc}
\usepackage[utf8]{inputenc}
\usepackage[russian]{babel}
\usepackage{tabularx}
\input{glyphtounicode}

\pagestyle{fancy}
\fancyhf{}
\fancyfoot{}
\renewcommand{\headrulewidth}{0pt}
\renewcommand{\footrulewidth}{0pt}

% Adjust margins
\addtolength{\oddsidemargin}{-0.5in}
\addtolength{\evensidemargin}{-0.5in}
\addtolength{\textwidth}{1in}
\addtolength{\topmargin}{-.5in}
\addtolength{\textheight}{1.0in}

\urlstyle{same}

\raggedbottom
\raggedright
\setlength{\tabcolsep}{0in}

% Sections formatting
\titleformat{\section}{
  \vspace{-4pt}\scshape\raggedright\large
}{}{0em}{}[\color{black}\titlerule \vspace{-5pt}]

\pdfgentounicode=1

%-------------------------
% Custom commands
\newcommand{\resumeItem}[1]{
  \item\small{
    {#1 \vspace{-2pt}}
  }
}

\newcommand{\resumeSubheading}[4]{
  \vspace{-2pt}\item
    \begin{tabular*}{0.97\textwidth}[t]{l@{\extracolsep{\fill}}r}
      \textbf{#1} & #2 \\
      \textit{\small#3} & \textit{\small #4} \\
    \end{tabular*}\vspace{-7pt}
}

\newcommand{\resumeProjectHeading}[2]{
    \item
    \begin{tabular*}{0.97\textwidth}{l@{\extracolsep{\fill}}r}
      \small#1 & #2 \\
    \end{tabular*}\vspace{-7pt}
}

\renewcommand\labelitemii{$\vcenter{\hbox{\tiny$\bullet$}}$}

\newcommand{\resumeSubHeadingListStart}{\begin{itemize}[leftmargin=0.15in, label={}]}
\newcommand{\resumeSubHeadingListEnd}{\end{itemize}}
\newcommand{\resumeItemListStart}{\begin{itemize}}
\newcommand{\resumeItemListEnd}{\end{itemize}\vspace{-5pt}}

%-------------------------------------------
\begin{document}

%----------HEADING----------
\begin{center}
    \textbf{\Huge \scshape Амангелды Шынгысхан} \\ \vspace{1pt}
    \small [Телефон] $|$ \href{mailto:email@example.com}{\underline{email@example.com}} $|$
    \href{https://linkedin.com/in/username}{\underline{linkedin.com/in/username}} $|$
    \href{https://github.com/username}{\underline{github.com/username}} $|$
    \href{https://t.me/username}{\underline{t.me/username}}
\end{center}


%-----------EDUCATION-----------
\section{Образование}
  \resumeSubHeadingListStart
    \resumeSubheading
      {Astana IT University}{Астана, Казахстан}
      {Бакалавр, 3-й (последний) курс}{2024 -- 2027 (ожидается)}
  \resumeSubHeadingListEnd


%-----------EXPERIENCE-----------
\section{Опыт работы}
  \resumeSubHeadingListStart

    \resumeSubheading
      {АО <<Казахстан Темир Жолы>>}{Астана, Казахстан}
      {DevOps and ML Engineer (стажировка)}{Янв. 2026 -- Июнь 2026}
      \resumeItemListStart
        \resumeItem{\textbf{Сервис оптимизации логистики:} в составе команды разработал систему для снижения 8--12-часового простоя локомотивов на сети из 1000+ станций для крупного акционерного общества (участника процесса IPO). Реализовал ML-модель для автоподбора оптимальных кандидатов к 90 локомотивам.}
        \resumeItem{\textbf{Геоданные \& Симуляция:} внедрил интерактивные карты на базе OpenStreetMap (OSM) и OpenRailwayMap. Спроектировал среду <<песочниц>> (sandbox) для симуляции и тестирования графиков движения и ситуаций ЧС.}
        \resumeItem{Проект не получил окончательного одобрения руководства; потенциал сокращения расходов оценивается в 40--50 тыс.\ тг/час (до 500 тыс.\ тг на локомотив).}
      \resumeItemListEnd

    \resumeSubheading
      {Розничный бизнес на Kaspi.kz}{Казахстан}
      {Главный технический специалист и аналитик продаж}{Июнь 2026 -- Сент. 2026}
      \resumeItemListStart
        \resumeItem{Настройка и обслуживание парка оборудования: моноблоки, роутеры, IP-камеры, ноутбуки, принтеры.}
        \resumeItem{Разработал внутренний сервис для сотрудников: автонакладные Kaspi, быстрая печать собственных этикеток и штрихкодов, уведомления о заказах Kaspi.}
        \resumeItem{Автоматизировал выгрузку каталога в YML для 2ГИС Услуги (4000+ товаров) — парсинг Kaspi через сессионные cookies, с уникальными ссылками на покупку и фотографиями.}
        \resumeItem{Поиск товаров по фото (база 14~000+ изображений, дообучение поверх готовой модели OpenAI) — сервисом ежедневно пользуются 25+ сотрудников.}
        \resumeItem{Автоматический сбор и анализ остатков склада Umag (12~360 SKU, 270~219 ед.) с выгрузкой метрик по нулевой себестоимости, отрицательным остаткам и заканчивающимся товарам отдельными отчётами.}
        \resumeItem{Автоматизация печати накладных Kaspi под термопринтеры; для бухгалтерии — автосоставление коммерческих предложений, ЭСФ и накладных.}
        \resumeItem{Анализ 4000+ активных товаров на Kaspi и Umag: сезонность, чистая выручка, исключение позиций без себестоимости.}
        \resumeItem{Извлечение 10~000+ уникальных номеров заказов Kaspi за последние 3 месяца для аналитики.}
        \resumeItem{Ведение отчётности по продажам и управление онлайн-базой склада.}
      \resumeItemListEnd

    \resumeSubheading
      {ТОО <<Wipon>>}{Казахстан}
      {Data Scientist and Machine Learning Engineer (стажировка)}{Сент. 2026 -- 15 Дек. 2026}

    \resumeSubheading
      {Tirek Systems}{Казахстан}
      {CTO, сооснователь}{Дек. 2024 -- настоящее время}
      \resumeItemListStart
        \resumeItem{Проектировал платформу для ускорения создания академических работ и анализа больших данных — \textbf{\href{https://perricheno.com}{perricheno.com}}. Участник Astana Hub Startups. <<Исследуй, пиши, анализируй, цитируй, визуализируй — всё в одном месте>>.}
        \resumeItem{\textbf{LLM Агенты:} проектировал асинхронную многослойную мультиагентную систему с послойными генерациями в 6 стадий. Встроил внутренний сервис компиляции LaTeX (на базе TeX Live Full), написанный на Go — ускорив процесс компиляции в 4 раза.}
        \resumeItem{\textbf{Data Analytics \& Компиляторы:} интегрировал изолированный сервис компиляции. Платформа анализирует миллионы строк датасетов; агенты пишут скрипты на R / Python / LaTeX / TikZ для статистического анализа и генерируют точные, редактируемые визуализации более 40+ видов.}
        \resumeItem{Настройка CI/CD-воркфлоу на GitLab с автоматическим тестированием.}
        \resumeItem{Инфраструктура на Kubernetes (k3s) и Linux (Ubuntu).}
      \resumeItemListEnd

  \resumeSubHeadingListEnd


%-----------PROJECTS-----------
\section{Проекты и научная деятельность}
    \resumeSubHeadingListStart
      \resumeProjectHeading
          {\textbf{Исследование E-commerce логистики} $|$ \emph{R, логистическая регрессия}}{Янв. -- Февр. 2026}
          \resumeItemListStart
            \resumeItem{В соавторстве с профессором Dr. Kamal Imran Mohd Sharif (PhD in Technology, Operation and Logistics, Universiti Utara Malaysia) провёл количественное исследование аналитики последней мили.}
            \resumeItem{Обработал массив из 10~000~000+ транзакций на R с применением бинарной логистической регрессии для оценки компромисса между скоростью и надёжностью доставки.}
          \resumeItemListEnd
      \resumeProjectHeading
          {\textbf{AI-First Financial Strategy Synthesis} $|$ \emph{Python, XGBoost, Gemini API}}{Февр. 2026}
          \resumeItemListStart
            \resumeItem{Соавтор: Yedil Talasbekov (ML Engineer) · курс AI for Finance Decisions (AIB 2402).}
            \resumeItem{\textbf{Генеративный AI-эксперимент:} 10-недельный итеративный цикл прогнозирования акций (GOOGL, ORCL, NET) на Gemini 3.0 Pro / 2.5 Pro / 2.0 Flash с фактчекингом по 10-K отчётам и live-поиском; средняя MAE снижена на 93\%, достигнув \$3.5 по Cloudflare и \$2.7 по Alphabet.}
            \resumeItem{\textbf{Ensemble-модель для прогнозирования:} архитектура на XGBoost, Random Forest и Logistic Regression — 36 предобученных моделей на все горизонты прогноза, R\textsuperscript{2} стабильно выше 0.98 на краткосрочных прогнозах.}
            \resumeItem{\textbf{Feature engineering:} 8-мерный вектор признаков (RSI, MACD, Bollinger Bands, SMA 5/20) для повышения качества сигнала.}
            \resumeItem{\textbf{MLOps и деплой:} полный CI/CD на Render.com и GitHub, асинхронная обработка live-данных через yfinance, автоматическая генерация графиков на matplotlib.}
          \resumeItemListEnd
    \resumeSubHeadingListEnd


%-----------TECHNICAL SKILLS-----------
\section{Технические навыки}
 \begin{itemize}[leftmargin=0.15in, label={}]
    \small{\item{
     \textbf{Языки программирования}{: Python, R} \\
     \textbf{Базы данных \& Аналитика}{: PostgreSQL, Pandas, Matplotlib, Seaborn, Power BI, Tableau} \\
     \textbf{Backend \& Архитектура}{: Микросервисы, Multi-Agent Systems, LLM, XGBoost} \\
     \textbf{DevOps \& Cloud}{: Linux (Ubuntu/Debian), Google Cloud Platform, Kubernetes (k3s), Docker, CI/CD (GitHub Actions, GitLab CI), Cloudflare} \\
     \textbf{Инструменты \& Геоданные}{: Git / GitHub / GitLab, OpenStreetMap (OSM), LaTeX, TikZ, Figma} \\
     \textbf{Языки}{: Русский (C1), Казахский (родной), Английский (B2 — технический)}
    }}
 \end{itemize}


%-------------------------------------------
\end{document}
